\scanpage{164}\mspage{109}
Donnons nous $S$ dans $\unsure{M|\mathcal M}$, considérons
\[
W_S\subset\mathrm{Fl}(M_{/S})
\]
l'ens de flèches induit par $W$ ($f:$
\diagram{\begin{tikzcd}[ampersand replacement=\&,column sep=small,row sep=tiny]
X \arrow[rr] \arrow[dr] \& \& Y \arrow[dl] \\
\& S \&
\end{tikzcd}}
est dans $W_S$ ssi il l'est en oubliant $S$), et désignons par $W_S$ (par abus de notation) l'ens de flèches induit par $W_S$ sur $\underline{\mathrm{Fib}}\,S$, ou plus généralement sur \deleted{\unsure{toute}} sous-catégorie pleine $\unsure{\mathcal F|F}$ de $M_{/S}$, avec
\[
\underline{\mathrm{Fib}}\,S\subset\mathcal F\subset M_{/S}.
\]

\textcircled{\unsure{I|II}} \deleted{\illegible{1}} Le foncteur
\[
(\underline{\mathrm{Fib}}\,S)W_S^{-1}\xrightarrow{\simeq}(M_{/S})W_S^{-1}\tag*{$(*)$}
\]
\inserted{induit par l'inclusion} est une équivalence de catégories, plus généralement pour \deleted{\unsure{toute}} $\mathcal F$ comme dessus
\[
(\underline{\mathrm{Fib}}\,S)W_S^{-1}\longrightarrow\mathcal F W_S^{-1}\longrightarrow(M_{/S})W_S^{-1}
\]
sont des \unsure{équiv.} de catégories. Un foncteur quasi-inverse est obtenu (dans les deux cas A) et B)) à l'aide du foncteur
\[
\gamma:M_{/S}\longrightarrow(\underline{\mathrm{Fib}}\,S)W_S^{-1}
\]
qui est donné sur les objets par
\[
\gamma(X)=C(X),
\]
et qui sur les flèches se définit de façon évidente dans le cas A (\unsure{fonctorialité|fonctorialité ?} de $C(X)$ en $X$), et dans le cas de B se construit ad-hoc, comme la seule loi fonctorielle pour laquelle les hom
\[
i_X:X\longrightarrow C(X)
\]
% archive sheet number: 163
